\section{Discussion}
\label{sec:discussion}

\subsection{Evaluation commitments}

\begin{itemize}
  \item \textbf{Cross-dataset evaluation.} Domain shift is the most likely
        cause of real-world failure; a random split of the training set
        flatters a model and predicts nothing about deployment.
  \item \textbf{Report the gap.} In-domain and out-of-domain figures appear
        side by side.
  \item \textbf{Event-level evaluation on sequences}, not frame-level on
        stills. The validator only demonstrates value over time.
  \item \textbf{$F_2$ rather than $F_1$ as the headline}, for the reason given
        in \cref{eq:fbeta-ratio}.
  \item \textbf{Never quote a figure we did not measure.} Published
        third-party numbers appear only in clearly attributed rows.
  \item \textbf{Every proportion carries its interval} (\cref{sec:intervals}).
\end{itemize}

\begin{recorded}{A trap we walked into and documented}
Our first false-alarm measurement on ignition sequences gave 0.8\,\%, against
71\,\% on curated hard negatives with the same model and threshold. The
difference is negative-set difficulty: frames shortly before a fire are
usually clear skies, whereas curated negatives are cloud, fog and dust chosen
to confuse. Had we published the 0.8\,\% figure as evidence the validator
works, any reviewer inspecting the negative set would have dismantled the
claim. Every false-alarm figure in this document names the negative set it
came from.
\end{recorded}

\subsection{Execution backends}
\label{sec:backends}

Two backends exist and are not interchangeable. \textsc{reference} runs the
as-published dynamic-shape graph on the CPU provider and is deterministic;
every figure in this document comes from it. \textsc{fast} runs a static-shape
variant on Apple's CoreML provider and is 1.7--1.9$\times$ faster
(\SI{34}{\milli\second} against \SI{60}{\milli\second} at \SI{1024}{\pixel}).

\textsc{fast} is not bit-identical. Across 100 validation frames the detection
count matched on 100 of 100 and boxes agreed to within \SI{0.41}{\pixel}, but
confidences drift by up to $0.012$ because the Neural Engine computes at lower
precision. At the $0.20$ operating point that drift changed no metric; at
$0.10$ and $0.30$ it moved box precision and recall by about one percentage
point by flipping a single detection across the threshold. \textsc{fast} is
therefore used for anything a person watches, and \textsc{reference} for
anything reported.

\begin{recorded}{A parity check that checked nothing}
The ONNX export parity check originally traced the model on random noise. A
trained detector finds nothing in random noise, so the check compared zero
detections against zero detections and passed unconditionally. It now traces a
real frame and reports \textsc{unverified} if that frame yields no detections,
instead of passing quietly.
\end{recorded}

\subsection{Findings}
\label{sec:findings}

Distilled from the measurements, phrased for reuse beyond this system:

\begin{enumerate}
  \item \textbf{Report suppression as two numbers or not at all.} Filtering
        and deduplication differ by a factor of eight between our hazards
        (\cref{tab:decomposition}); one figure conflates a false-alarm claim
        with alert housekeeping.
  \item \textbf{Run the threshold-matched control \emph{and} the transfer
        experiment before claiming alarm reduction.} On negatives it had seen,
        an oracle threshold beat our cascade; across unseen cameras that same
        threshold either let false alarms drift up or lost a detection, while
        the untuned cascade held both. The oracle result alone would have
        misled in either direction; both controls cost an afternoon on cached
        detections.
  \item \textbf{Persistence has a price list:}
        $\mathbb{E}[\Delta t] = (N-1)/f$. The latency cost of temporal
        confirmation is a property of the camera's sampling rate, computable
        before any data is collected.
  \item \textbf{Pooled metrics flatter minority-critical classes.} Selecting
        the aerial flood model on pooled IoU would have chosen a model
        better at recognising dry ground (\cref{sec:pooling}); stratify on
        the case the system exists for.
  \item \textbf{Appearance priors fail where the confuser shares the
        property} — fog is as grey as smoke — and geometric priors can
        invert safety on the severe tail (the worst floods are the largest).
        Measure each level of a life-safety cascade on the negatives it will
        actually face, and be willing to delete.
  \item \textbf{Image-level metrics can hide a dead detector.} A label-shift
        bug survived image-level evaluation at precision $0.83$ and was
        exposed only at box level, at recall $0.002$
        (\cref{sec:twolevels}). Evaluate at both levels, always.
  \item \textbf{Interval or it didn't happen.} At $n = 12$ images, a recall
        of $1.000$ means $[0.757, 1.000]$ (\cref{tab:intervals}); reporting
        the interval is the difference between a measurement and an
        anecdote.
\end{enumerate}

We would argue the transferable contribution is the practice: every
proportion carrying the interval its sample supports, stratified metrics
where pooling flatters, suppression decomposed into the claim it supports
and the housekeeping it is not, controls run against one's own headline, and
twelve retracted claims documented in place. None of this required novel
architecture. All of it changed what the paper was allowed to say.

\subsection{Limitations and outstanding work}
\label{sec:outstanding}

Listed in the order it would be tackled. The first three are shortages of
available material rather than defects in the system.

\begin{enumerate}
  \item \textbf{No temporal hard-negative set exists.} Tuned thresholding
        suffices on easy negatives (\cref{sec:threshold-matched}); no
        threshold survives the curated hard stills; whether persistence
        survives them is untestable without hard negatives that are also
        sequences — fog rolling over a ridge, filmed. Assembling one from
        HPWREN's archive is the single most valuable experiment this project
        has not run.
  \item \textbf{Building access rests on \BuildingTestImages{} held-out
        images.} Strong on that split (civilian box recall
        \BuildingBoxRecall, interval \BuildingCivilianRecallCI{} over
        instances), but twelve positive images cannot make the image-level
        figures precise (\cref{tab:intervals}), and the released data cannot
        be widened. It needs a different post-disaster aerial dataset with
        people annotated.
  \item \textbf{No urban flood footage.} The flood demonstration uses a
        public-domain gauge camera on a creek. The 2024 Valencia DANA is the
        event this system is motivated by, and no footage of a flooded
        street is yet available to it under a licence permitting publication.
  \item \textbf{The road-incident detector has no standalone measurement}
        against held-out labels; it remains the validator's generality
        experiment only.
  \item \textbf{The gauge correlation is one event on one camera}
        ($r = +0.50$, \cref{sec:gage}); it should be repeated across many
        USGS cameras and events before the trend claim is leaned on.
  \item \textbf{The aerial flood segmenter has no held-out split}
        (\cref{tab:flood-aerial}); FloodNet's withheld masks make this
        unfixable from the public release.
  \item \textbf{Night operation.} Fire and flood are evaluated in daylight
        only; the flood segmenter mislabels infrared frames and the fetch
        pipeline excludes them by envelope. Disasters do not keep daylight
        hours — the Valencia flooding peaked in the evening — so an
        infrared-domain evaluation, likely fine-tuned on the USGS cameras'
        own night frames, is required before any claim extends past dusk.
\end{enumerate}
